\documentclass[a4paper]{article}
\usepackage{amsfonts}
\usepackage{amsmath}
\usepackage{mathtools}
\usepackage{amssymb}
\usepackage{multicol}
\usepackage{fullpage}
\usepackage{graphicx}
\usepackage{polynom}
\title{Facultad de Matemática, Astronomía, Física y Computación}
\author{Didáctica}
\date{11 de Agosto de 2026}
\begin{document}
\maketitle
\section{Resolucion}
    \begin{enumerate}
        \item Tarea 1
        \begin{enumerate}
            \item Basta con entender que por el teorema fundamental de la aritmetica, todo numero entero puede ser expresado como un producto de numeros primos, y que este producto es unico salvo el orden de los factores. Por lo tanto, es múltiplo de 5 y 35 purs en los factores que componen la multiplicación están el 15 y 14 que ambos son el producto de 5 y 7 pero el producto de 5 y 7 es 35, aunque los multiplos se dan en el producto con numeros naturales, particularmente nuestra operacion cuenta con la suma de dos numeros positivos, lo que resulta imposible forzar que se haga negativa al multiplicarla por mas naturales.
            \item numero entero de k tal que $5^2(k+3)^2$ sea multiplo de 7, como 5 es un numero primo y estoy viendo el producto, basta con forzar que $(k+3)^2$ sea multiplo de 7, por lo tanto $k+3$ es multiplo de 7, entonces $k$ es multiplo de 7 menos 3, es decir $k=7n-3$ para algun numero entero n. Y si, K puede ser negativo porque al estar elevado al cuadrado el resultado es positivo, y por lo tanto el producto es multiplo de 7.
        \end{enumerate}
        \item 
    \end{enumerate}
74
\end{document}